% SOURCE_AUTHORITY: sga5_fr_workpass.tex, lines 10281--10300 (LF-normalized unit).
% SOURCE_SHA256: 791F4EFFC5E02832D5D77ED03518C8156D6F07E4C8238B03545DB93D883FBB28
\begin{statement}{Proposição 6.2}
(Sra. RAYNAUD [5].) Sejam $S$ um esquema e $p:G\to S$ um $S$-esquema em grupos redutivo (SGA 3 XIX). Para todo $A_S$-módulo localmente constante construtível $M$ e para todo inteiro $i$, o $A_S$-módulo $\R^ip_*(p^*M)$ é localmente constante e construtível. Além disso, para todo morfismo $u:S'\to S$ que define um diagrama cartesiano
\[
\begin{tikzcd}[column sep=large,row sep=large]
G' \arrow[r,"u'"] \arrow[d,"p'"'] & G \arrow[d,"p"]\\
S' \arrow[r,"u"] & S,
\end{tikzcd}
\]
o morfismo de mudança de base
\[
u^*\R^ip_*(p^*M)\to \R^ip'_*(p\circ u')^*(M)
\]
é um isomorfismo.
\end{statement}

\emph{Demonstração.} As duas afirmações do enunciado são locais para a topologia étale de $S$. Por (SGA 3 XII 2.3), pode-se, portanto, supor que $G$ é cindido. Seja então $B$ um subgrupo de Borel de $G$ do tipo indicado em (SGA 3 XXII 5.1.1), de modo que o quociente $G/B$ existe por (SGA 3 XXII 5.8.5), e sejam $q:G\to G/B$ e $r:G/B\to S$ os morfismos canônicos. Como o morfismo $r$ é próprio e liso, satisfaz o teorema de mudança de base próprio, e os funtores $\R^ir_*$ transformam feixes abelianos localmente constantes construtíveis em feixes do mesmo tipo (SGA 4 XVI 2.2). Por outro lado, segundo (SGA 3 XXII 5.5.1), o $S$-esquema $B$ é isomorfo a um produto finito de fibrados vetoriais e de fibrados vetoriais privados da seção zero; portanto, de (1.1) e (1.2) deduz-se imediatamente que, denotando por $t:B\to S$ o morfismo canônico, os feixes $\R^it_*(t^*M)$, onde $M$ é um $A_S$-módulo localmente constante construtível, são localmente constantes construtíveis e estáveis por mudança de base. Seja agora $M$ um $A_S$-módulo localmente constante construtível. Das considerações precedentes e de (6.1) deduz-se que os feixes $\R^iq_*(p^*M)$ são localmente constantes construtíveis; utilizando a sequência espectral de Leray
\[
\R^ir_*(\R^jq_*)(p^*M)\Rightarrow \R^{i+j}p_*(p^*M),
\]
conclui-se que o mesmo vale para os feixes $\R^ip_*(p^*M)$. Finalmente, por (6.1), os feixes $\R^iq_*(p^*M)$ são estáveis por mudança de base; graças ao teorema de mudança de base próprio para $r$, deduz-se que os feixes $\R^ip_*(p^*M)$ também são estáveis por mudança de base.
